% Generated draft: requires mathematical and source review.
\documentclass[11pt]{article}
\usepackage[margin=25mm]{geometry}
\usepackage{fontspec}
\setmainfont{Latin Modern Roman}
\usepackage{amsmath,amssymb,mathtools,unicode-math}
\setmathfont{Latin Modern Math}
\usepackage{enumitem,needspace,xurl,hyperref}
\setlength{\parindent}{0pt}
\setlength{\parskip}{5pt}
\setlength{\emergencystretch}{3em}
\newcommand{\N}{\mathbb{N}}
\newcommand{\Z}{\mathbb{Z}}
\newcommand{\Q}{\mathbb{Q}}
\newcommand{\R}{\mathbb{R}}
\newcommand{\C}{\mathbb{C}}
\newcommand{\F}{\mathbb{F}}
\newcommand{\A}{\mathbb{A}}
\newcommand{\PP}{\mathbb{P}}
\newcommand{\E}{\mathbb{E}}
\newcommand{\eps}{\varepsilon}
\newcommand{\dd}{\,\mathrm{d}}
\begin{document}
\section*{UnsolvedMath ID: 800266}
\textbf{Source problem code:} catalogue entry
\subsection{Definitions and mathematical statement}
Let \(f:\C^2\to\C^2\) be a polynomial automorphism, meaning a polynomial map with a polynomial inverse. Assume, as in the source, that \(f\) is not conjugate to an elementary map. In the usual triangular convention, an elementary map has the form
\[
(z,w)\longmapsto (az+p(w),bw+c),\qquad a,b\in\C\setminus\{0\},\quad c\in\C,\quad p\in\C[w].
\]
The source does not specify the category of conjugacy; this hypothesis retains that unresolved qualification.

Define the forward locally bounded set by
\[
\mathcal F_b(f)=\left\{x\in\C^2:\ \text{there exist an open neighborhood }U\ni x\text{ and }M<\infty\text{ such that }\|f^j(y)\|\le M\text{ for every }y\in U\text{ and }j\ge1\right\}.
\]
The source calls this the Fatou set. Let \(\Omega\) be a connected component of \(\mathcal F_b(f)\), and suppose that \(f^n(\Omega)=\Omega\) for some positive integer \(n\). Taking \(n>0\) is equivalent to the source's condition that a nonzero integer period exists, since \(f\) is invertible. Put \(G=f^n\). All boundaries below are taken in \(\C^2\).

The source calls \(\Omega\) recurrent if it contains a point whose orbit does not converge to \(\partial\Omega\), and nonrecurrent otherwise. It does not specify the meaning of convergence to the boundary as a set or whether this criterion uses the full orbit under \(f\) or the return orbit under \(G\). These qualifications remain unresolved here.

For recurrent \(\Omega\), the source asserts the alternatives: (a) a basin of attraction of a sink orbit, (b) a basin of attraction of a Siegel disk, or (c) a basin of a Herman ring. A sink orbit is a periodic orbit whose return derivative has all eigenvalues of modulus less than one. The terms Siegel disk and Herman ring refer to one-dimensional invariant rotation domains modeled, respectively, on a disk and an annulus, with return dynamics holomorphically conjugate to an irrational rotation. The source does not specify the relevant return iterate or the precise meaning of their basins in \(\C^2\). It reports that examples of (a) and (b) are readily constructed.

\begin{enumerate}
\item[(i)] Can alternative (c), a recurrent periodic component that is the basin of a Herman ring, occur under these hypotheses? The precise basin condition requires clarification from the source.
\item[(ii)] For every nonrecurrent periodic component \(\Omega\) as above, does there exist a single point \(P\in\partial\Omega\) such that
\[
\forall Q\in\Omega,\qquad \lim_{k\to\infty}G^k(Q)=P?
\]
Here the limit is ordinary convergence in \(\C^2\), and the point \(P\) must work for every \(Q\in\Omega\).
\item[(iii)] For every such nonrecurrent \(\Omega\), does there exist \(P\in\partial\Omega\) with \(G(P)=P\) and a root-of-unity eigenvalue of the derivative denoted by \(DF(P)\) in the source? The source does not define \(F\). If its intended meaning is \(F=f^n=G\), the resulting conditional formulation is
\[
\exists P\in\partial\Omega\ \exists p\in\Z\ \exists q\in\Z_{>0}:\qquad G(P)=P,\qquad e^{2\pi i p/q}\in\operatorname{Spec}(DG(P)).
\]
This identification of \(F\) is an unresolved interpretation, not an additional supplied hypothesis.
\end{enumerate}

The source's statement that essentially nothing was known about nonrecurrent periodic domains is historical commentary from the 1995 list, not a verified assessment of subsequent literature.
\subsection{Short English statement}
For a polynomial automorphism of \(\C^2\) outside the elementary class, consider periodic components where forward iterates are locally bounded. The questions ask whether a recurrent component can be a Herman-ring basin, and whether every nonrecurrent component has a common boundary limit under the return map and a boundary periodic point with a root-of-unity multiplier. The source leaves the conjugacy category, recurrence and basin conventions, and the derivative in the last question insufficiently specified.
\subsection{Sources}
\begin{itemize}
\item Local UnsolvedMath registry, permanent ID \(800266\), exact source problem code \texttt{catalogue entry}; title: \emph{Several Complex Variables Problem List --- Problem 65}. The supplied registry labels the record \emph{open}; that label is not independently verified here.
\item Buzzard, Fornaess, and Krantz, \emph{The Several Complex Variables Problem List} (1995), Problem 65. Supplied source URL: \url{https://arxiv.org/abs/math/9509202}. The supplied provenance records access on 2026-10-16 and extraction from arXiv LaTeX; no independent source visit is claimed.
\item The individual problem is signed E. Bedford, with the supplied contact addresses \texttt{bedford@iubacs.bitnet} and \texttt{bedford@ucs.indiana.edu}. The registry lists Buzzard, Fornaess, and Krantz as proposers and 1995 as the proposed year; the individual attribution to Bedford is retained separately.
\end{itemize}
\end{document}
